% Abstract (<= 250 words; count with: detex sections/abstract.tex | wc -w)
\begin{abstract}
A UAV-borne 5G base station can serve search-and-rescue robots, but it must be
managed autonomously, without knowing the terrain. We
built a terrain-blind, learning-based management layer for such a cell, which
forecasts link failure and places the UAV in three dimensions, and flew it in
closed loop on an OpenAirInterface testbed serving four robots.
Validated on held-out terrain and against null controls, the largest gain
came not from learning but from fixing a silent failure of the managed system. Long forest outages turned
out to be unbounded uplink queuing inside the user equipment, not coverage
loss: constant-rate video exceeded weak links, and command-and-control (C2)
traffic waited minutes behind it. A per-robot uplink cap driven only by
round-trip time and bytes sent removed these outages: in a forest patrol it raised C2 availability from
66.8\,\% to 91.8\,\% and cut the time with no robot connected from 17.4\,\% to
0.9\,\%, and it improved every placement policy. After censoring a second
silent failure, dead UE processes that contaminated 12.4\,\% of observations, a linear
model forecasts C2 failure 20\,s ahead with modest skill (ROC-AUC 0.667) that
falls as model capacity grows, and it does not recover how terrain changes the
effect of altitude. In closed loop, no learned placement beat fixed hover or the
robot centroid; an exploring altitude learner lost 18--27 points of the time
with all robots connected, and a cautious one was not distinguishable from them. Bounding
queues and checking liveness should precede learning where to fly.
\end{abstract}
% Introduction (1.0 p). Claims: paper/CLAIMS.md; novelty: paper/NOVELTY.md
\section{Introduction}
\label{sec:intro}

After an earthquake, flood or fire, the ground robots sent to search a
disaster site need a mission-critical link for command and control (C2) and
video, and the terrestrial network they would use is often gone or
overloaded~\cite{murphy2014disaster}. A 5G base station (gNB) carried by an
unmanned aerial vehicle (UAV) can bring a private cell to the
robots~\cite{zeng2016wireless,mozaffari2019tutorial}, but someone, or
something, must manage it: decide where to fly it, notice when a robot is
about to lose its link, and keep the links usable under the load the robots
put on them. On site there is no operator to tune it and no map of a terrain
that may have just been changed, so this management has to be autonomous,
proactive and blind to the terrain.

Our earlier testbed study~\cite{makhubela2026satnac} showed that terrain
decides whether moving such a node helps: a roaming node raised a weak robot's
C2 availability in open field and cost about 22 points in forest. A natural next
step is a learning-based management layer that predicts link failures and
places the UAV in three dimensions, without being told the terrain. This paper
builds that layer, flies it in closed loop over a real 5G stack, and asks what
it can and cannot deliver. The answer turned out to depend less on the learning
than on whether the managed system itself was measured and loaded soundly.

We test four claims, each against an explicit control (Section~\ref{sec:setup}):
\begin{itemize}
\item[\textbf{C1}] Long forest outages under robot video are caused by
unbounded uplink queuing inside the UE, not by coverage, and a delay-adaptive
uplink cap at each robot removes them without starving strong links.
\item[\textbf{C2}] A terrain-blind model forecasts the onset of C2 failure
20\,s ahead on an unseen terrain better than chance, and higher-capacity models
do worse.
\item[\textbf{C3}] The same model does not recover how terrain changes the
effect of altitude, so it cannot choose altitude.
\item[\textbf{C4}] Learned placement does not beat fixed hover or the robot
centroid; an exploring altitude learner loses to them and a cautious one does
not.
\end{itemize}
C1 holds; C2 and C4 hold weakly; C3 holds as a negative result.

\textit{Contributions.} The work falls under wireless and 5G networks,
mission-critical networks, proactive and autonomic management, AI-based
management, monitoring and measurement, and testbed experimentation. Beyond
the testbed and terrain comparison of~\cite{makhubela2026satnac}, it
contributes:
\begin{enumerate}
\item \emph{A diagnosis and an autonomic fix.} We show that the forest
outages are a silent queueing failure of the managed system. A per-robot cap
driven only by round-trip time and bytes sent, with no change to the radio
access network, raised forest C2 availability from 66.8\,\% to 91.8\,\% and improved each of
four placement policies, one run each (Section~\ref{sec:r-cap}).
\item \emph{A measurement-integrity detector.} A second silent failure, the
UE process dying after radio-link failure, contaminated 12.4\,\% of training
observations; we give a mechanistic rule that removes it
(Section~\ref{sec:dead}).
\item \emph{An evaluation of terrain-blind learning for proactive management},
on held-out terrain and in closed loop, with the negative results stated:
modest forecasting skill, worse with capacity, and no recovery of the
terrain-dependent altitude effect (Sections~\ref{sec:r-forecast}
and~\ref{sec:r-flip}).
\item \emph{A closed-loop account of learned placement}, showing why a single
mission is too short to learn the forest's altitude effect and how a cautious
learner avoids the resulting loss (Section~\ref{sec:r-place}).
\end{enumerate}
All results come from emulation over real 5G software; a hardware trial is the
next step.
% Related work (0.5 p). Refs: SATNAC list (author-supplied) + web-verified (refs.bib verify_note)
\section{Background and Related Work}
\label{sec:related}

\textit{Aerial base stations and their testbeds.} The elevation-dependent
line-of-sight model~\cite{alhourani2014modeling} underlies most work on where
to place a UAV base station, from 3-D placement~\cite{boryaliniz2016efficient,
kalantari2016number} and trajectory design~\cite{wu2018common} to deep
reinforcement learning~\cite{qiu2020placement}; surveys cover the wider
field~\cite{zeng2016wireless,mozaffari2019tutorial}. These works optimise
coverage or rate under an analytic or simulated channel. Testbeds now run
OpenAirInterface on UAVs with O-RAN control~\cite{moore2024aerpaw} or emulate
5G mobility with ray tracing~\cite{iye2025owdt}. Our earlier
study~\cite{makhubela2026satnac} coupled a real 5G stack to a terrain
simulator and found that terrain decides whether moving the node helps.

\textit{Learning-based and autonomous management.} Zero-touch management aims
at fully autonomous operation under human oversight~\cite{coronado2022zsm}.
Machine learning for link-quality estimation~\cite{cerar2021lqe}, fault
management~\cite{mulvey2019fault} and failure prediction in 5G RANs, with deep
and transformer models~\cite{islam2023dnn,hasan2024gentrap}, learns from the
history of a deployed network. A UAV arriving at a disaster site has no such
history for its terrain, so we test on an unseen terrain and in closed loop.

\textit{Buffering in cellular links.} Buffers that hold rather than drop
packets inflate delay~\cite{gettys2012bufferbloat}, and per-user cellular
buffers are a known source~\cite{jiang2012tackling}. Remedies range from active
queue management~\cite{nichols2012codel} and delay-aware rate
control~\cite{winstein2013sprout} to, recently, RAN-side congestion marking:
in 5G, time spent in the RLC buffer dominates one-way
delay~\cite{wan2025l4span}, and retransmission and uplink-scheduling buffers can
add delay comparable to congestion~\cite{liu2025fog}. Those remedies change
network queues, the RAN or the transport protocol. Our cap applies the classic AIMD
rule~\cite{chiu1989aimd} to a delay signal at the robot, because constant-rate
robot video does not back off and the UE's buffer is not ours to change.
% System and method. Sources: docs/SYSTEM_ARCHITECTURE.md, scripts/ml/{onset_forecast,
% planner,online_planner,altitude_bandit,geo_features}.py, scripts/adaptive_uplink_cap.py
\section{System and Method}
\label{sec:method}

\subsection{Platform}
\label{sec:platform}
The managed system is the protocol-in-the-loop testbed
of~\cite{makhubela2026satnac}, extended from two to four robots
(Fig.~\ref{fig:loop}). An OpenAirInterface (OAI)~\cite{nikaein2014oai} gNB and
four OAI UEs run over the OAI radio simulator and register to an
Open5GS~\cite{open5gs} core. Each UE's TUN interface sits in its own Linux
network namespace, so robot traffic crosses the real 5G user plane. Gazebo
moves the UAV and four ground robots; a bridge reads their poses every 0.5\,s
and writes a terrain-specific air-to-ground path loss into each UE's channel:
free-space loss with an elevation-dependent line-of-sight
term~\cite{alhourani2014modeling}, plus foliage~\cite{weissberger1982},
diffraction~\cite{itup526} or building loss by terrain, and correlated
shadowing~\cite{gudmundson1991} (details in~\cite{makhubela2026satnac}). A
watchdog relaunches a UE whose process exits after radio-link failure.

\begin{figure}[t]
\centering
\includegraphics[width=0.92\columnwidth]{fig_loop.pdf}
\caption{The closed management loop. Shaded boxes are the management
functions studied here: a per-robot uplink cap and a terrain-blind
forecaster, planner and altitude learner.}
\label{fig:loop}
\end{figure}

\subsection{Measurement Integrity: Dead-Process Detection}
\label{sec:dead}
In OAI, a radio-link failure terminates the UE process. Before the watchdog
existed, the robot's C2 then stayed down for the rest of the run, and the gNB's
per-UE counters froze. Those cycles are labelled as link failures, but they
record a dead process, and a model with lagged features detects them trivially
(the first difference of the block error rate is exactly zero). We flag a cycle
as dead when, for at least three consecutive cycles, the gNB counters
\texttt{dl\_bler}, \texttt{ul\_nb\_rbs} and \texttt{ul\_bits\_per\_sym} are
bit-identical to the previous cycle \emph{and} C2 is down. The conjunction
matters: frozen counters on a working link are a quiet cycle, and a failed
cycle with moving counters is a real degradation. The rule flags 12.4\,\% of
all cycles and 72.2\,\% of hills failure labels (Table~\ref{tab:a1}); flagged
cycles are censored from training and evaluation.
% A4: dead_process_extent.csv; rule: onset_forecast.py dead_ue_mask (FREEZE_MIN)

\subsection{Terrain-Blind Learning Layer}
\label{sec:layer}
\textit{Forecaster.} The label is the onset of C2 failure: for a robot whose C2
is up at time $t$, is it down at $t+20$\,s? C2 is up when fewer than 50\,\% of
its probes are lost. Features are UAV--robot geometry, gNB-side link
adaptation and error statistics, and the last commanded move, each with lags,
differences and rolling statistics within one robot's series (65 features).
Three inputs are forbidden: the simulator's injected path loss, the computed
received power (algebraically equivalent to it), and any terrain label. A guard
function checks every feature matrix before it reaches an estimator. Models are
trained leave-one-terrain-out, with thresholds and calibration chosen on an
inner time split of the training terrains.

\textit{Planner.} Each 30-s decision, the planner scores candidate UAV
positions (absolute altitudes 25--55\,m in 5\,m steps, crossed with lateral
offsets of 0 or 15\,m) by a soft minimum over the robots' predicted risks,
$J(u) = \beta^{-1}\log\sum_i e^{\beta r_i(u)}$ with $\beta = 10$, so the
worst-served robot dominates. It moves toward the minimiser at most 10\,m in
altitude per decision. The planner's risk model is the same logistic
regression with spline terms on altitude, distances and elevation angle and
Platt calibration (pooled held-out ROC-AUC 0.659).

\textit{Altitude learner.} Because the model cannot choose altitude
(Section~\ref{sec:r-flip}), altitude is learned during the mission. The learner
holds an incumbent altitude and rewards each altitude by the fraction of robots
with C2 up. Its estimates are time-blocked (240\,s blocks), so only comparisons
made close in time count. Every few decisions it probes a neighbouring altitude
($\pm$5 or 10\,m), and it moves when the lower confidence bound of the
neighbour-minus-incumbent difference is above zero ($z = 0.5$). The
\emph{cautious} variant adds a home altitude, the hover altitude of 40\,m: a
move further from home must clear 1.5 standard errors \emph{and} a 0.02 gain,
a move back toward home only the normal bar, and probes are rarer (every 4th
decision, every 20th once all neighbours are worse).
% planner.py (Z_TARGETS, DXY_R, MAX_DZ_STEP, beta=10); altitude_bandit.py AltitudeProber; online_planner.py CAUTIOUS

\subsection{Delay-Adaptive Uplink Cap}
\label{sec:cap}
The cap moves the bottleneck from the UE into the robot's kernel, where it can
be bounded and prioritised, and lets its rate follow the link. It needs no
change to the RAN or core. On each robot's TUN interface, an HTB shaper has a
ceiling equal to the cap $c$. C2 and ICMP use a high-priority class with a
20-packet queue; video uses a low-priority class with a 10-packet queue, so
excess video is dropped rather than buffered. Each robot pings the core every
0.2\,s. ICMP bypasses the kernel video queue, so its RTT rises only when the
UE's own buffer backs up. With $r_0$ the minimum RTT over 60\,s and $\tilde r$
the median over the last second, Algorithm~\ref{alg:cap} runs every 0.2\,s.
The ``in use'' test keeps an idle robot's cap from rising and then overloading
the link when video resumes. The controller re-creates the shaper when a UE
reattaches with a new TUN. Its only inputs, RTT and bytes sent, are available
on real hardware.

\begin{algorithm}[t]
\caption{Per-Robot Uplink Cap (Every 0.2\,s)}
\label{alg:cap}
\begin{algorithmic}[1]
\small
\STATE initially $c \leftarrow 200$\,kb/s, $\mathit{slow} \leftarrow$ true
\IF{$\tilde r > r_0 + 150$\,ms \OR no reply for 1\,s}
  \STATE $c \leftarrow \max(50,\,0.7c)$; $\mathit{slow} \leftarrow$ false; hold 2\,s
\ELSIF{$\tilde r < r_0 + 50$\,ms \AND sent $\ge 0.7c$ \AND 2\,s since last raise}
  \STATE $c \leftarrow \min(1600,\ \mathit{slow}\ ?\ 2c : c + 50)$
\ENDIF
\end{algorithmic}
\end{algorithm}
% adaptive_uplink_cap.py constants: LO,HI,START=50,1600,200; CUT .7; STEP 50; HOLD 2; UP_EVERY 2; HIGH 150; LOW 50; BASE_WIN 60
% Experimental setup. Sources: EVIDENCE.md A-E, run_*_flyoff.sh, video_streams.py
\section{Experimental Setup}
\label{sec:setup}

\textit{Traffic and metrics.} Each robot sends constant-rate UDP video at
1.5\,Mb/s to the core and is probed for C2 by ping; a robot's C2 is up when
fewer than 50\,\% of its probes are lost. These definitions differ from those
of~\cite{makhubela2026satnac} (two robots, 8\,Mb/s video bursts, C2 up when
any of five echoes returns), so numbers are not comparable across the two
papers. Placement runs are reported as the time-weighted share of a mission
with mean C2 up, with all four robots up, and with no robot up, and the
90th-percentile video latency. Forecasts are scored by PR-AUC relative to the
no-skill rate (lift) and by ROC-AUC, pooled over held-out terrains and on the
worst held-out terrain.

\textit{Data.} The forecasting and altitude-effect results use 54,572
robot-cycles from 20 missions of 90--200\,min over four terrains (open field,
urban, forest, hills), all flown before the cap existed. The cap was tested in
a static load test in forest and open field (6\,min of video between 2-min
no-load phases), in a 30-min forest patrol, and in the placement fly-offs.

\textit{Placement arms.} Every arm flies the same pilot, seed and 90-min
mission and differs only in how x--y and altitude are chosen
(Table~\ref{tab:arms}): the pilot decides every 30\,s, changes altitude by at
most 10\,m per decision, and stays within a 20--60\,m geofence. Arm B, the
robot centroid at fixed altitude, is the baseline of our research proposal; arm
A, fixed hover, is the simplest policy; the repositioning policy
of~\cite{makhubela2026satnac} is a further baseline for the risk planner.

\begin{table}[t]
\caption{Placement Arms}
\label{tab:arms}
\centering
\small
% Table 1 (hand-written: definitions, no numbers). Source: run_xyz_flyoff.sh,
% scripts/ml/online_planner.py (SAR_PLANNER_XY / _ALT / _ALT_LEARNER).
\begin{tabular}{@{}llll@{}}
\toprule
Arm & x--y & Altitude & Role \\
\midrule
A & held & fixed 40\,m & baseline (hover) \\
B & robot centroid & fixed 40\,m & baseline (proposal) \\
C & robot centroid & learned on the spot & learned altitude \\
D & risk model & learned on the spot & learned x--y--z \\
\bottomrule
\end{tabular}

\end{table}

\textit{Validation protocol.} A result is accepted only after it survives the
control that matches its threat: leave-one-terrain-out evaluation with
worst-fold reporting for forecasting; censoring of dead-process cycles for every
label; a same-policy null control (the same comparison applied to two nights of
one policy) for closed-loop comparisons made across nights; and the uplink cap
for any forest outcome. Confidence intervals are 95\,\% block bootstraps over
contiguous stretches (10-min blocks for missions), because consecutive cycles
are strongly correlated; an effect whose interval includes zero is reported as
not established. Per-arm placement comparisons use one 90-min run per arm and
within-run intervals, which do not capture run-to-run variation.
\todo{update: three seeds per arm in forest and open field from the repeat campaign}
% Results. Every number: see paper/figs/README.md and paper/EVIDENCE.md (IDs in comments).
\section{Results}
\label{sec:results}

\subsection{C1: Uplink Queuing, Not Coverage, and an Autonomic Cap}
\label{sec:r-cap}

\textit{The long forest outages are a queue inside the UE.} With the four
robots parked under canopy (24--30\,dB uplink attenuation) and the UAV at
40\,m, C2 was clean without load. Once the four 1.5\,Mb/s video feeds started,
uplink ping RTT grew from about 60\,ms to minutes within the first seconds and
kept rising for the whole load period, peaking at 350--444\,s
(Fig.~\ref{fig:rtt}(a)); 64--91\,\% of pings were lost and, for 64\,\% of the
time, all eight probes were lost at once (Table~\ref{tab:static}). Probes were
delayed rather than dropped, and the links drained only 0.13--0.49\,Mb/s of the
1.5\,Mb/s offered. The kernel queue on each TUN held no backlog and dropped
nothing, so the excess, about 79\,MB per robot, was accepted and held by the OAI
UE (E1). After the video stopped, C2 stayed dead for 14 to more than 18\,min,
consistent with each link slowly draining its backlog.

% E1: tab_static_load (analyse.py), ul_*.txt max RTT, CONGESTION_TEST_REPORT.md (qdisc, 79 MB)
\begin{figure*}[t]
\centering
\includegraphics[width=\textwidth]{fig_rtt.pdf}
\caption{Uplink C2 ping RTT per robot in the static forest load test (log
scale; video on in the shaded interval). (a) Without a cap, RTT grows to
minutes; traces end where replies stop before the test ends. (b) With the
adaptive cap, RTT stays near its unloaded level.}
\label{fig:rtt}
\end{figure*}

\begin{table}[t]
\caption{Static Forest Load Test, Four Robots, 6\,min of Video}
\label{tab:static}
\centering
\footnotesize\setlength{\tabcolsep}{3pt}
% generated by tab_static_load.py -- do not edit
\begin{tabular}{@{}lccc@{}}
\toprule
 & No cap & Fixed 100\,kb/s & Adaptive \\
\midrule
Uplink ping loss, video on & 64--91\,\% & 0\,\% & 0\,\% \\
Uplink ping RTT, median & 3--4\,min & 69--87\,ms & 76--108\,ms \\
All probes lost (on / off) & 64 / 100\,\% & 0 / 0\,\% & 0 / 0\,\% \\
Recovery after video & 14--18+\,min & none & none \\
Video goodput (Mb/s) & 0.13--0.49 & 0.10 & 0.10--0.42 \\
Video latency, median & 2.2--2.8\,min & 1.08--1.11\,s & 0.30--1.08\,s \\
\bottomrule
\end{tabular}

\end{table}

\textit{Bounding the queue removes the outages.} A fixed 100\,kb/s cap below
every link's drain rate, applied at the same placement, removed all C2 loss and
the recovery tail, but discarded most of the video (Table~\ref{tab:static}).
The adaptive cap of Section~\ref{sec:cap} also gave zero C2 loss while settling
at a different rate on each link (medians 105--429\,kb/s, 32--40 cuts per robot
in 6\,min), delivering 0.10--0.42\,Mb/s of video at a median latency of
0.30--0.42\,s on three robots; the fourth, with the weakest link, ran at
1.08\,s, the time its 10-packet video queue takes to drain at about 100\,kb/s
(Fig.~\ref{fig:rtt}(b)). In open field, with 2--9\,dB of attenuation, every cap
rose to its 1.6\,Mb/s ceiling, made no cuts, and carried the full 1.55\,Mb/s
per robot at 82--89\,ms median latency (E4): in this test, where links were
strong, the cap stayed out of the way.
% E2, E3: cap_log.csv (medians, cuts), tab_static_load; E4: CONGESTION_TEST_REPORT follow-up 3

\textit{The cap lifts every placement policy.} In a 30-min forest patrol with
the UAV hovering at 40\,m, the cap raised mean C2 availability from 66.8\,\% to
91.8\,\% against the same seed uncapped, cut the time with no robot connected
from 17.4\,\% to 0.9\,\%, and the 90th-percentile video latency from 73\,s to
1.6\,s (Table~\ref{tab:a5}). Across the four placement arms of the forest
fly-off (Fig.~\ref{fig:caparms}), the time with no robot connected fell from
25--35\,\% to 0--0.3\,\% and mean C2 rose from 51--67\,\% to 82--89\,\% in
every arm. Each comparison is one run per arm on different nights, but the
direction is the same in all four.
% E5: tab_a5_patrol; E6: fig_cap_arms (summary CSV, sets xyz/capped)

\begin{figure}[t]
\centering
\includegraphics[width=\columnwidth]{fig_cap_arms.pdf}
\caption{Forest fly-off, each placement arm without and with the adaptive cap
(one 90-min run per bar; whiskers: 95\,\% block-bootstrap CI). Arms C and D
used a UCB altitude bandit uncapped and the prober learner capped.}
\label{fig:caparms}
\end{figure}

The same mechanism may bear on the forest result of our earlier
study~\cite{makhubela2026satnac}, where a roaming node drove 8\,Mb/s video
bursts to over 70\,\% loss; its traffic pattern differed, and we have not
tested this.

\subsection{C2: Terrain-Blind Forecasting Has Modest Skill, and Capacity Hurts}
\label{sec:r-forecast}

\textit{A linear model forecasts failure onset 20\,s ahead on unseen terrain,
weakly.} After dead-process censoring (Section~\ref{sec:dead}), the best
model, an L1-regularised logistic regression on 65 terrain-blind features,
reaches a pooled PR-AUC 1.79 times the no-skill rate and ROC-AUC 0.667, and no
held-out terrain falls below ROC-AUC 0.538 (Fig.~\ref{fig:capacity}). Skill
peaks at a 20-s horizon: a six-feature trend model scores ROC-AUC 0.650 at
20\,s and 0.552 at 60\,s (Table~\ref{tab:a2}). Forest labels in this data
predate the cap and are partly queue backlog (Section~\ref{sec:r-cap}).
% A6: phase2c_model_class.csv; A5: phase2b_onset.csv

\begin{figure}[t]
\centering
\includegraphics[width=\columnwidth]{fig_capacity.pdf}
\caption{Forecasting skill by model class: pooled PR-AUC lift over no skill
on held-out terrains, and the worst held-out-terrain ROC-AUC (65 features,
20-s onset horizon, 38,303 onset observations, leave-one-terrain-out).}
\label{fig:capacity}
\end{figure}

\textit{More capacity is worse.} On identical folds and features the ordering
is monotone in capacity: three linear models lead, tree ensembles follow, and
three multilayer perceptrons trail. The largest network scores ROC-AUC 0.503,
and 0.404 on its worst terrain, below chance. One exception should be stated:
on six trend features at a 60-s horizon, the medium network beats the linear
models on pooled lift (1.55 against 1.21--1.22) but not on its worst terrain (ROC
0.456 against 0.550). Sequence models on 10\,Hz radio traces add nothing: on
3,115 labelled cycles over three folds, a TCN and a GRU average ROC-AUC 0.493
and 0.523 against 0.540 for a logistic regression on window summaries (A7).
Under leave-one-terrain-out validation with four terrains, we take worst-fold
skill, not pooled skill, as the figure of merit; our reading is that it rewards
models that learn a direction rather than thresholds tied to the training
terrains.
% A6 exception: phase2c n==38803 block, trend6 mlp_medium 1.551 / worst 0.456 vs logistic_L2 1.221 / 0.551; A7: phase6_seq.csv

\subsection{C3: The Model Does Not Recover How Terrain Changes the Altitude Effect}
\label{sec:r-flip}

The measured effect of climbing differs strongly between terrains: per +10\,m,
C2 availability gains 0.208 more in open field than in forest per cycle, and
0.167 more on a uniform 60-s time base that removes the outcome-dependent
cycle length (Table~\ref{tab:contrasts}). This terrain dependence was reported
in~\cite{makhubela2026satnac}; the question here is whether a terrain-blind
model recovers it. It does not. The same contrast computed from the model's
risk estimates is $+0.006$ ($p = 0.365$) per cycle and $+0.010$ on the 60-s
base, both indistinguishable from zero; at 120\,s it reaches $+0.021$, which is
significant but about one sixth of the measured $+0.117$. Two of the other
three contrasts are significantly inverted per cycle.
% B1-B4: tab_contrasts (phase4b/4c/4d); 120 s: phase4d_interactions recovered_raw

\begin{table}[t]
\caption{Altitude Effect Contrasts (C2 Availability per +10\,m): Measured
Versus Recovered From the Model. $^*$: 95\,\% CI Excludes 0}
\label{tab:contrasts}
\centering
\small
\begin{tabular}{@{}c@{}}
{\footnotesize per cycle \hspace{12mm} 60-s time base}\\[-1pt]
% generated by tab_contrasts.py -- do not edit
\begin{tabular}{@{}lcccc@{}}
\toprule
Contrast & Cycle & 60-s bins & $h{=}20$\,s & $h{=}60$\,s \\
\midrule
Open $-$ forest & $+0.208$$^*$ & $+0.167$$^*$ & $+0.001$ & $+0.006$ \\
Open $-$ hills & $+0.153$$^*$ & $+0.138$$^*$ & $-0.003$ & $+0.037$$^*$ \\
Urban $-$ forest & $+0.073$$^*$ & $+0.070$$^*$ & $+0.006$ & $-0.020$$^*$ \\
Hills $-$ forest & $+0.055$$^*$ & $+0.030$ & $+0.003$ & $-0.031$$^*$ \\
\bottomrule
\end{tabular}
% measured: per cycle and 60-s uniform bins; recovered (LightGBM): per cycle at horizons 20 s and 60 s

\end{tabular}
\end{table}

The consequence for control is direct. Replayed through the planner on each
held-out terrain, no surrogate chooses altitude sensibly: linear models choose
about 55\,m in every terrain, and a gradient-boosted model chooses about 47\,m in both
forest and open field (D1). Altitude therefore has to be learned during the
mission.
% D1: logs/ml/altsel_run.log

\subsection{C4: Learned Placement Under the Cap}
\label{sec:r-place}

\textit{An exploring altitude learner loses to the baselines.} With the cap on,
hover (A) and centroid (B) kept all four robots connected 70.5\,\% and
67.0\,\% of the time in forest, indistinguishable from each other (B$-$A
$-4.0$ [$-20.4$, $+11.2$] points). The two arms with an exploring learner were
worse: C$-$A $-18.1$ [$-34.1$, $-2.6$] and D$-$A $-27.4$ [$-43.0$, $-10.7$].
They lost by climbing to 45--55\,m (Fig.~\ref{fig:alt}).
% D5: capped summary + CAPPED_FOREST_FLYOFF_REPORT bootstrap table

\textit{One mission cannot determine the sign of the forest altitude effect.}
Comparing arm C with arm B at matched times, with the same x--y rule and robot
paths, gives $-6.6$ points of C2 per +10\,m ($[-14.9, +0.9]$): probably harmful,
not established. The slope the learner itself can estimate within one 90-min
run is $+4.6$ ($[-9.1, +12.0]$) in arm C and $+5.9$ ($[-10.7, +14.9]$) in arm
D. Its interval spans both signs, so the learner moved on noise.
% D6: logs/ml/forest_altitude_effect_capped.csv

\begin{figure*}[t]
\centering
\includegraphics[width=\textwidth]{fig_altitude.pdf}
\caption{UAV altitude during the capped forest missions for the two
learned-altitude arms. The exploring learner drifts to 45--55\,m; the cautious
learner keeps 40\,m as its incumbent and leaves it only for short probes.}
\label{fig:alt}
\end{figure*}

\textit{A cautious learner avoids that loss.} Homed at 40\,m and requiring 1.5 standard
errors plus 0.02 of reward to move away, the learner kept 40\,m as its
incumbent in both arms and spent 83--84\,\% of each mission there. All four
robots were connected 65.2\,\% (C) and 72.0\,\% (D) of the time, not
distinguishable from the baselines, though the intervals are wide (C$-$A
$-5.1$ [$-17.9$, $+6.8$]; D$-$A $+1.6$ [$-13.4$, $+17.7$]),
and arm D gained $+28.8$ [$+11.9$, $+45.1$] points over its exploring version.
D's earlier deficit therefore came from altitude, not from the model's x--y.
No learned arm beat a baseline. \todo{replace with the 3-seed forest and
open-field results of the repeat campaign (Fig.~\ref{fig:seeds}); decision rule
fixed in CLAIMS.md}
% D7: cautious_forest_flyoff summary; time at 40 m: flyoff_forest_cautious_analyse.py

\begin{figure}[t]
\centering
\includegraphics[width=\columnwidth]{fig_arms_seeds.pdf}
\caption{All four robots C2 up per arm with the cap, cautious learner for C
and D; bar = mean over runs, dots = individual 90-min runs. \todo{partial:
forest seed 1 only; regenerate after the repeat campaign}}
\label{fig:seeds}
\end{figure}

\textit{The risk planner survives its null control in two terrains only.}
Against the repositioning policy of~\cite{makhubela2026satnac}, matched on
geometry, the planner gained $+10.5$ [$+7.0$, $+13.2$] points in open field and
lost $-15.3$ [$-24.8$, $-5.0$] in forest; the same comparison between two
nights of the \emph{same} baseline policy gave $-0.5$ and $-5.0$, both not
significant. In urban and hills that null was as large as the effect ($+18.7$
and $-21.7$), so those comparisons are uninterpretable
(Table~\ref{tab:a3}).
% C2: planner_vs_twophase_nullctl.csv
% Limitations (0.4 p); full version in Appendix C. Sources: docs/CHANNEL_MODEL_LIMITATIONS.md, EVIDENCE.md
\section{Limitations and Threats to Validity}
\label{sec:limits}

The radio is real 5G software, but propagation is injected: the channel has
large-scale loss and correlated shadowing but no multipath, terminal height is
fixed at 0.5\,m, and absolute availability is internally consistent rather than
externally calibrated. The OAI UE may buffer uplink data differently from a
commercial modem, and constant-rate video is a worst case for the queue;
whether C1 transfers is the first question for hardware. Every
between-terrain claim rests on four terrains. The forecasting and
altitude-effect data predate the cap, so their forest labels are partly queue
backlog. Placement comparisons use one 90-min run per arm with within-run
intervals, and the robots' safety guards react to link state, so the same seed
does not give identical patrols across arms. Appendix~\ref{app:threats} gives
the full list.
\todo{revise the last-but-one sentence once the three-seed repeats are in}
% Conclusion (0.35 p)
\section{Conclusion}
\label{sec:conclusion}

We built a terrain-blind, learning-based management layer for a UAV-borne 5G
cell serving four search-and-rescue robots and flew it in closed loop over real
5G software. The most effective management action was not learned: a per-robot
uplink cap driven by round-trip time removed a silent queueing failure inside
the UE and improved each of four placement policies, one run each (C1). The learning layer forecasts
link-failure onset 20\,s ahead with modest skill that falls as model capacity
grows (C2), does not recover how terrain changes the effect of altitude (C3),
and, in closed loop, does not beat fixed hover or the robot centroid; made
cautious, it is not distinguishable from them within wide intervals (C4). For autonomous management of such
cells, the lesson is to verify that the managed system is measured and loaded
soundly, bounding its queues and checking that its processes are alive, before
learning where to fly it. The next steps are a hardware trial of the cap and
retraining the layer on capped, watchdog-protected data.
